\subsection{Module 1: cAMP Receptor Binding and Desensitisation}\label{sec:m1}

\paragraph{Role.}
cAMP is sensed by the G-protein-coupled receptor cAR1.
Binding to intact aggregation-competent cells is heterogeneous: \mk{M1-slow4}{about 4\,\% of the sites dissociate slowly ($K_d=12.5$\,nM)}~\cite{vanhaastert1984}, \mk{M1-fast-high}{about 40\,\% are fast, high-affinity sites ($K_d=60$\,nM)}~\cite{vanhaastert1984} and \mk{M1-fast-low}{about 60\,\% are fast, low-affinity sites ($K_d=450$\,nM)}~\cite{vanhaastert1984}; \mk{M1-conv9}{under cAMP the high-affinity sites convert to low-affinity ones with first-order kinetics ($t_{1/2}\approx9$\,s)}~\cite{vanhaastert1984}.
Occupancy \mk{M1-ser-phos}{induces phosphorylation of serine clusters in the C-terminal domain}~\cite{vaughan1988,hereld1994}, which is \mk{M1-llb-corr}{correlated with the loss of ligand binding}~\cite{caterina1995a,caterina1995b} and \mk{M1-aff-red}{reduces the affinity of the receptor rather than removing it from the surface}~\cite{caterina1995a,caterina1995b}.
The model resolves the two fast classes, and represents the low-affinity class as the phosphorylated receptor; the slow class (4\,\% of sites) is omitted.
Each receptor form has the binding topology of the \mk{M1-bis-red}{reduced-order receptor model} of Biswas et al.~\cite{biswas2021} (Table~1, rows 27/28), $\mathrm R+\mathrm{cAMP}\rightleftharpoons\mathrm{RC}$, to which the model adds a phosphorylated pair; the \mk{M1-bis-full}{full three-class model} of~\cite{biswas2021} is available as an option (Table~\ref{tab:m1brx}) and is not used here.

\begin{table*}[!t]
\centering
\caption{Module~1 reactions (default two-state receptor). R is the unphosphorylated cAR1 receptor, R$_p$ its phosphorylated form, RC and R$_p$C the cAMP-bound states; A is the cytosolic adapter of reactions 1.8. A coloured reaction is the one to which the equally coloured statement in the text and the equally coloured passage in the cited paper refer; a small square appends a further statement.}
\label{tab:m1rxn}
\footnotesize
\renewcommand{\arraystretch}{1.15}
\begin{tabularx}{\textwidth}{@{}l>{\raggedright\arraybackslash}p{0.37\textwidth}>{\raggedright\arraybackslash}X@{}}
\toprule
ID & Reaction & Propensity (rate law)\\
\midrule
\mkn{M1-bis-red}{M1-kon-tab,M1-fast-high}{1.1} & \mkn{M1-bis-red}{M1-kon-tab,M1-fast-high}{$\mathrm{R}+\mathrm{cAMP}_e \to \mathrm{RC}$} & $k_\mathrm{on}[\mathrm{R}][\mathrm{cAMP}_e]$ \\
\mkn{M1-bis-red}{M1-kon-tab,M1-fast-high,M1-koff1s}{1.2} & \mkn{M1-bis-red}{M1-kon-tab,M1-fast-high,M1-koff1s}{$\mathrm{RC} \to \mathrm{R}+\mathrm{cAMP}_e$} & $k_\mathrm{off}[\mathrm{RC}]$ \\
\mkn{M1-fast-low}{M1-aff-red,M1-aff35}{1.3} & \mkn{M1-fast-low}{M1-aff-red,M1-aff35}{$\mathrm{R}_p+\mathrm{cAMP}_e \to \mathrm{R}_p\mathrm{C}$} & $k_{\mathrm{on},p}[\mathrm{R}_p][\mathrm{cAMP}_e]$ \\
\mkn{M1-fast-low}{M1-aff-red,M1-aff35,M1-koff1s}{1.4} & \mkn{M1-fast-low}{M1-aff-red,M1-aff35,M1-koff1s}{$\mathrm{R}_p\mathrm{C} \to \mathrm{R}_p+\mathrm{cAMP}_e$} & $k_{\mathrm{off},p}[\mathrm{R}_p\mathrm{C}]$ \\
\mkn{M1-phos45}{M1-ser-phos,M1-llb-corr,M1-conv9}{1.5} & \mkn{M1-phos45}{M1-ser-phos,M1-llb-corr,M1-conv9}{$\mathrm{RC} \to \mathrm{R}_p\mathrm{C}$} & $k_\mathrm{ph}[\mathrm{RC}]$ \\
\mkt{M1-basal}{1.5b} & \mkt{M1-basal}{$\mathrm{R} \to \mathrm{R}_p$} & $k_{\mathrm{ph}0}[\mathrm{R}]$ \\
\mkn{M1-loss80}{M1-h10}{1.6} & \mkn{M1-loss80}{M1-h10}{$\mathrm{R}_p\mathrm{C} \to \mathrm{RC}$} & $k_\mathrm{dp}[\mathrm{R}_p\mathrm{C}]$ \\
\mkt{M1-deph2}{1.7} & \mkt{M1-deph2}{$\mathrm{R}_p \to \mathrm{R}$} & $k_\mathrm{b}[\mathrm{R}_p]$ \\
\mkt{M1-fcd}{1.8a} & \mkt{M1-fcd}{$\mathrm{RC}+\mathrm{A} \to \mathrm{RC}+\mathrm{A}^*$} & $k_{A}[\mathrm{RC}][\mathrm{A}]$ \\
\mkt{M1-fcd}{1.8b} & \mkt{M1-fcd}{$\mathrm{A}^* \to \mathrm{A}$} & $k_{A,\mathrm{off}}[\mathrm{A}^*]$ \\
\bottomrule
\end{tabularx}
\end{table*}

\begin{table*}[!t]
\centering
\caption{Module~1 constants (default receptor model).}
\label{tab:m1const}
\footnotesize
\renewcommand{\arraystretch}{1.15}
\begin{tabularx}{\textwidth}{@{}l r l c >{\raggedright\arraybackslash}X@{}}
\toprule
Symbol & Value & Unit & Basis & Meaning, justification, source\\
\midrule
$R_\mathrm{tot}$ & \num{3077} & $\mathrm{molec/voxel}$ & E & 40\,001 receptors per cell ($0.64\,\mu$M). Cells developed for 4\,h carry about $7\times10^4$ binding sites~\cite{johnson1991}, and Biswas et al.\ use $7\times10^4\pm5\times10^3$~\cite{biswas2021}; the model value is a lower design choice. \\
$k_\mathrm{on}$ & \num{7.5} & $\mu\mathrm{M}^{-1}\mathrm{s}^{-1}$ & L & Fast high-affinity class H of Ref.~\cite{biswas2021}, Table~1, taken from~\cite{vanhaastert1984}. \\
$k_\mathrm{off}$ & \num{0.45} & $\mathrm{s^{-1}}$ & L & Gives $K_d=k_\mathrm{off}/k_\mathrm{on}=60$\,nM, the fast high-affinity class ($\approx\!40$\,\% of sites) of~\cite{vanhaastert1984}. \\
$\varepsilon_p$ & \num{0.125} &  & C & Ratio of on-rates of R$_p$ and R. Chosen so that $K_{d,p}=480$\,nM, close to the 450\,nM low-affinity class of~\cite{vanhaastert1984}; the loss of affinity on phosphorylation is documented in~\cite{caterina1995b,caterina1995a}. \\
$k_{\mathrm{on},p}$ & \num{0.938} & $\mu\mathrm{M}^{-1}\mathrm{s}^{-1}$ & D & $\varepsilon_p k_\mathrm{on}$. The affinity loss is carried by the on-rate, so R$_p$ still binds and releases cAMP. \\
$k_{\mathrm{off},p}$ & \num{0.45} & $\mathrm{s^{-1}}$ & L & Equal to $k_\mathrm{off}$; the fast dissociation half-life of $\approx\!1$\,s is common to both fast classes~\cite{vanhaastert1984}. \\
$k_\mathrm{ph}$ & \num{0.02} & $\mathrm{s^{-1}}$ & L & Agonist-induced phosphorylation ($t_{1/2}=35$\,s) against $t_{1/2}\approx45$\,s at saturating cAMP in~\cite{vaughan1988}. \\
$k_{\mathrm{ph}0}$ & \num{5e-4} & $\mathrm{s^{-1}}$ & E & Ligand-independent phosphorylation (40$\times$ slower); a basal component exists~\cite{vaughan1988,hereld1994}. \\
$k_\mathrm{dp}$ & \num{5e-3} & $\mathrm{s^{-1}}$ & C & Dephosphorylation of the bound receptor ($t_{1/2}=139$\,s). Sets the phosphorylated fraction $k_\mathrm{ph}/(k_\mathrm{ph}+k_\mathrm{dp})=80\,\%$ under sustained saturating cAMP, matching the $>\!80\,\%$ loss of high-affinity binding~\cite{vanhaastert1984,johnson1991}. \\
$k_\mathrm{b}$ & \num{5.78e-3} & $\mathrm{s^{-1}}$ & L & $t_{1/2}=2$\,min, the receptor dephosphorylation half-time after removal of the stimulus~\cite{vaughan1988}. \\
$k_A,\ k_{A,\mathrm{off}}$ & \num{0.07},\ \num{0.02} & $\mu\mathrm{M}^{-1}\mathrm{s}^{-1},\ \mathrm{s^{-1}}$ & E & Adapter charged by RC ($\tau=50$\,s); dormant in the default network (see text). \\
$[\mathrm{A}]_\mathrm{tot}$ & \num{154} & $\mathrm{molec/voxel}$ & E & Adapter pool. \\
$D_{\mathrm{R}}$ & \num{0.024} & $\mu\mathrm{m^2/s}$ & L & Common lateral diffusion constant of transmembrane proteins in \emph{Dictyostelium}, $0.024\pm0.004\,\mu\mathrm{m}^2/\mathrm{s}$ ($n=27$)~\cite{takebayashi2023}; Biswas et al.\ use $0.027$~\cite{biswas2021,ueda2001}. \\
$D_{\mathrm{A}}$ & \num{10} & $\mu\mathrm{m^2/s}$ & E & Generic cytosolic value. \\
\bottomrule
\end{tabularx}
\end{table*}

\begin{table*}[!t]
\centering
\caption{Optional full multi-affinity receptor module of Biswas et al.~\cite{biswas2021} (Table~1, rows 1--26), selectable in the code but \emph{not} used in the results of this paper. Classes H, L, S: high, low and slow affinity; P is the free-phosphate pool; $^P$ marks phosphorylated states. Each row lists the forward and reverse constants.}
\label{tab:m1brx}
\footnotesize
\renewcommand{\arraystretch}{1.15}
\begin{tabularx}{\textwidth}{@{}l>{\raggedright\arraybackslash}p{0.37\textwidth}>{\raggedright\arraybackslash}X@{}}
\toprule
ID & Reaction & Propensity (rate law)\\
\midrule
1/2 & $\mathrm{H}+\mathrm{cAMP}\rightleftharpoons \mathrm{H{:}C}$ & $k_H=\num{7.5},\ k_{-H}=\num{0.45}$ \\
3/4 & $\mathrm{L}+\mathrm{cAMP}\rightleftharpoons \mathrm{L{:}C}$ & $k_L=\num{2.2},\ k_{-L}=\num{1}$ \\
5/6 & $\mathrm{S}+\mathrm{cAMP}\rightleftharpoons \mathrm{S{:}C}$ & $k_S=\num{4},\ k_{-S}=\num{0.05}$ \\
7/8 & $\mathrm{H}\rightleftharpoons \mathrm{L}$ & $k_{HL}=\num{0.08},\ k_{-HL}=\num{0.0536}$ \\
9/10 & $\mathrm{H{:}C}\rightleftharpoons \mathrm{L{:}C}$ & $k_{HLC}=\num{0.08},\ k_{-HLC}=\num{7.1e-3}$ \\
11/12 & $\mathrm{H{:}C}+\mathrm{P}\rightleftharpoons {}^P\mathrm{H{:}C}$ & $k_{HCP}=\num{5.3e-4},\ k_{-HCP}=\num{8e-4}$ \\
13/14 & ${}^P\mathrm{H}+\mathrm{cAMP}\rightleftharpoons {}^P\mathrm{H{:}C}$ & $k_{PHC}=\num{0.04},\ k_{-PHC}=\num{2.42e-3}$ \\
15/16 & $\mathrm{H}+\mathrm{P}\rightleftharpoons {}^P\mathrm{H}$ & $k_{PH}=\num{5.3e-4},\ k_{-PH}=\num{8e-4}$ \\
17/18 & $\mathrm{L{:}C}+\mathrm{P}\rightleftharpoons {}^P\mathrm{L{:}C}$ & $k_{LCP}=\num{5.3e-4},\ k_{-LCP}=\num{8e-4}$ \\
19/20 & ${}^P\mathrm{L}+\mathrm{cAMP}\rightleftharpoons {}^P\mathrm{L{:}C}$ & $k_{PLC}=\num{5},\ k_{-PLC}=\num{1.5}$ \\
21/22 & $\mathrm{L}+\mathrm{P}\rightleftharpoons {}^P\mathrm{L}$ & $k_{PL}=\num{5.3e-4},\ k_{-PL}=\num{8e-4}$ \\
23/24 & ${}^P\mathrm{H}\rightleftharpoons {}^P\mathrm{L}$ & $k_{PHL}=\num{4.3e-4},\ k_{-PHL}=\num{2.9e-4}$ \\
25/26 & ${}^P\mathrm{H{:}C}\rightleftharpoons {}^P\mathrm{L{:}C}$ & $k_{PHLC}=\num{4.3e-4},\ k_{-PHLC}=\num{3.8e-5}$ \\
\bottomrule
\end{tabularx}
\end{table*}


\paragraph{Reactions.}
\begin{description}\itemsep0pt
\item[1.1/1.2] Reversible cAMP binding to the high-affinity receptor R, with $K_d=60$\,nM. This is the fastest step of the cascade ($\tau=1.1$\,s at $c=K_d$) and is close to quasi-equilibrium with everything downstream.
\item[1.3/1.4] Binding to the phosphorylated receptor R$_p$. Affinity is lost through the on-rate ($K_{d,p}=480$\,nM), so R$_p$ still binds and releases cAMP, i.e.\ it is slower to engage, not locked out. A low-affinity receptor is still on the steep part of its binding curve at concentrations where the high-affinity receptor is saturated; R$_p$ therefore extends the dynamic range of the stage.
\item[1.5/1.5b] Phosphorylation, agonist-induced (1.5) and basal (1.5b, 40$\times$ slower). Agonist-induced phosphorylation is measured with \mk{M1-phos45}{$t_{1/2}\approx45$\,s at saturating cAMP}~\cite{vaughan1988} and \mk{M1-basal}{a ligand-independent basal component exists}~\cite{vaughan1988}; the model uses $t_{1/2}=35$\,s.
\item[1.6/1.7] Dephosphorylation of the bound (1.6) and free (1.7) receptor. The free receptor is dephosphorylated with \mk{M1-deph2}{$t_{1/2}=2$\,min after removal of the stimulus}~\cite{vaughan1988}, so 1.7 sets the time to resensitise after a pulse. Under sustained saturating cAMP, 80\,\% of receptors are phosphorylated at steady state ($k_\mathrm{ph}/(k_\mathrm{ph}+k_\mathrm{dp})$), matching \mk{M1-loss80}{the loss of more than 80\,\% of the surface binding sites during prolonged stimulation}~\cite{johnson1991} and \mk{M1-h10}{the fall of the high-affinity fraction from about 40\,\% to about 10\,\%}~\cite{vanhaastert1984}.
\item[1.8a/b] A cytosolic adapter A is charged by occupied receptors and relaxes with $\tau=50$\,s. It was introduced as a fold-change-detecting gate for a direct receptor-to-cyclase edge (reaction (8.11)), motivated by \mk{M1-fcd}{the fold-change response of the cAMP relay}~\cite{kamino2017}. That edge is switched off ($k=0$), so the adapter has \emph{no downstream effect} in the default network; it is retained as an observable.
\end{description}

\paragraph{Constants.}
$k_\mathrm{on}$ and $k_\mathrm{off}$ are \mk{M1-kon-tab}{the high-affinity constants of Table~1 of}~\cite{biswas2021}, \mk{M1-kon-tab}{which in turn come from the analysis of}~\cite{vanhaastert1984}, and give $K_d=60$\,nM.
The on-rate ratio $\varepsilon_p=1/8$ makes $K_{d,p}=480$\,nM, close to the 450\,nM low-affinity class; the phosphorylation-induced affinity change was reported as \mk{M1-aff35}{3--5-fold for one low-affinity class} in the literature used by~\cite{biswas2021}, so the 8-fold ratio is at the upper end and \mk{M1-ratio75}{equals the ratio of the two fast classes}~\cite{vanhaastert1984}.
The receptor number, 40\,001 per cell, is lower than the \mk{M1-rtot}{$\approx7\times10^4$ sites measured in cells developed for 4\,h}~\cite{johnson1991} and used in~\cite{biswas2021}; it is a design choice and not a measurement.
The lateral diffusion constant $0.024\,\mu\mathrm{m^2/s}$ is \mk{M1-dif-t}{the mean over 27 transmembrane proteins measured in \emph{Dictyostelium}}~\cite{takebayashi2023}, close to \mk{M1-dif-u}{the value $0.027$ used by Biswas et al.}~\cite{biswas2021,ueda2001}.
A much faster phospho-cycle ($k_\mathrm{ph}=1$\,s$^{-1}$, about 50 times faster than the measured half-time of 45\,s) was tested as a way to de-saturate the receptor stage and rejected, because it made the unphosphorylated receptor resupply-limited and hence no longer a sensor; the present values are the literature-anchored ones.

\paragraph{Limitations.}
Ligand-induced phosphorylation is not necessary for adaptation of the downstream responses or for chemotaxis; \mk{M1-kim-ok}{phosphorylation-deficient receptors still adapt}~\cite{kim1997}.
The phospho-cycle here therefore models only the loss of binding affinity, not the termination of signalling.
Identifying the low-affinity class with the phosphorylated receptor is a simplification, since the affinity conversion (\mk{M1-conv9}{$t_{1/2}\approx9$\,s}~\cite{vanhaastert1984}) is faster than phosphorylation (\mk{M1-phos45}{$t_{1/2}\approx45$\,s}~\cite{vaughan1988}).
Receptor internalisation is not modelled, consistent with the observation that \mk{M1-aff-red}{cAR1 remains at the surface}~\cite{caterina1995b}.
